% =====================================================================
% Anexo A. Repositorio de código (R2, 04/11/2026; revisado el 08/11/2026)
% Revisados el 08/11/2026: el token de ESIOS no está en los archivos
% versionados ni en el historial, ni hay rutas personales en ellos, y el
% historial tiene una sola identidad de autor. Licencia MIT.
% Pendiente antes de la entrega: hacer público el repositorio y crear la
% etiqueta v1.0 sobre el commit definitivo.
% =====================================================================
\chapter{Repositorio de código y reproducibilidad}\label{anexo:repositorio}

El código del trabajo, el protocolo de comparación y las tablas de resultados
que se presentan en la memoria están en un repositorio de GitHub:

\begin{center}
	\url{https://github.com/Guillermosoriano20/tfm-prevision-demanda-electrica}
\end{center}

\paragraph{Contenido.} El repositorio contiene los \textit{scripts} que
descargan y construyen el conjunto de datos, los que ajustan y evalúan los
cuatro modelos y los que calculan la capa probabilística y el desglose de los
resultados; el protocolo de comparación, con las decisiones de cada fase
fechadas antes de ejecutarla; y los archivos CSV de resultados en los que se
apoyan las tablas y las figuras de los capítulos~\ref{cap:metodologia}
y~\ref{cap:resultados}. Las semillas están fijadas en el código y las versiones
de las librerías, en el archivo de dependencias, con lo que se persigue que los
resultados puedan reproducirse.

\paragraph{Lo que no contiene.} No se publican los datos de entrada ni los
resultados intermedios, que ocupan decenas de megabytes y se regeneran con los
\textit{scripts} a partir de las fuentes (sección~\ref{sec:datos}). La descarga
de la demanda de Red Eléctrica exige un \textit{token} personal gratuito
de la API de e·sios, que se guarda en un archivo local de variables de entorno,
excluido del repositorio. La huella de contenido del conjunto
(tabla~\ref{tab:control}) permite comprobar que el conjunto regenerado es
idéntico al empleado.

\paragraph{Licencia.} El repositorio se distribuye con la licencia MIT (archivo
\texttt{LICENSE}). El archivo \texttt{README.md} explica el entorno, el orden de
los \textit{scripts} y las garantías de la evaluación final. Los datos de
entrada se rigen por las condiciones de uso de sus fuentes.

\paragraph{Versión citada.} La memoria corresponde a la versión del repositorio
etiquetada como \texttt{v1.0}.
